% GENERATED FILE. Edit canonical lessons, then run npm run generate:books.
% canonical-source-hash: fnv1a64:c2b4cedb69c73d0b
% canonical-lessons: TE-C161-natyam-ceyu, TE-C161-jarupuko, TE-C161-palakarincu, TE-C161-siggupadu, TE-C161-ascaryapovu

\chapter{To dance, To celebrate, To greet, To feel shy, To be surprised}
\label{ch:te-to-dance-to-celebrate-to-greet-to-feel-shy-to-be-surprised}
\hlchaptermodality{161}

\begin{chapteropening}
\textbf{By the end of this chapter:} \emph{I can say to dance, to celebrate, to greet, to feel shy and to be surprised.}

Complete the last lesson of chapter 161: I can say to dance, to celebrate, to greet, to feel shy and to be surprised.
\end{chapteropening}

\section[\texorpdfstring{\te{నాట్యం} \te{చేయు} (nāṭyaṁ cēyu)}{నాట్యం చేయు (nāṭyaṁ cēyu)}]{\te{నాట్యం} \te{చేయు} (nāṭyaṁ cēyu) — to dance}

\label{lesson:TE-C161-natyam-ceyu}



\begin{quote}
Before the new one: say the Telugu for to remember, then the Telugu for to agree.
\end{quote}

\subsection*{You'll want to know: \te{నాట్యం} \te{చేయు}}
\textbf{\te{నాట్యం} \te{చేయు}} — \emph{nāṭyaṁ cēyu} — \textquotedblleft{}to dance\textquotedblright{}.

\begin{grammarlens}[title={What you've built}]
One more word for this chapter.
\end{grammarlens}

\subsection*{Your turn}
\begin{itemize}
\raggedright
  \item \emph{Say it:} \emph{nāṭyaṁ cēyu}
  \item \emph{Say it:} \emph{nāṭyaṁ cēyu}, once more
  \item \emph{Say it:} say \emph{nāṭyaṁ cēyu}
  \item \emph{Recall:} say the Telugu for to remember, then the Telugu for to agree, then say \emph{nāṭyaṁ cēyu} again
\end{itemize}

\begin{tcolorbox}[breakable,colback=teal!4,colframe=teal!35!black,title={Before you move on}]
What does \te{నాట్యం} \te{చేయు} mean? (\textquotedblleft{}To dance\textquotedblright{}.) Say it once more.
\end{tcolorbox}
\section[\texorpdfstring{\te{జరుపుకో} (jarupukō)}{జరుపుకో (jarupukō)}]{\te{జరుపుకో} (jarupukō) — to celebrate (a festival or occasion)}

\label{lesson:TE-C161-jarupuko}



\begin{quote}
Before the new one: say the Telugu for to agree, then the Telugu for to dance.
\end{quote}

\subsection*{You'll want to know: \te{జరుపుకో}}
\textbf{\te{జరుపుకో}} — \emph{jarupukō} — \textquotedblleft{}to celebrate (a festival or occasion)\textquotedblright{}.

\begin{grammarlens}[title={What you've built}]
One more word for this chapter.
\end{grammarlens}

\subsection*{Your turn}
\begin{itemize}
\raggedright
  \item \emph{Say it:} \emph{jarupukō}
  \item \emph{Say it:} \emph{jarupukō}, once more
  \item \emph{Say it:} say \emph{jarupukō}
  \item \emph{Recall:} say the Telugu for to agree, then the Telugu for to dance, then say \emph{jarupukō} again
\end{itemize}

\begin{tcolorbox}[breakable,colback=teal!4,colframe=teal!35!black,title={Before you move on}]
What does \te{జరుపుకో} mean? (\textquotedblleft{}To celebrate (a festival or occasion)\textquotedblright{}.) Say it once more.
\end{tcolorbox}
\section[\texorpdfstring{\te{పలకరించు} (palakariñcu)}{పలకరించు (palakariñcu)}]{\te{పలకరించు} (palakariñcu) — to greet}

\label{lesson:TE-C161-palakarincu}



\begin{quote}
Before the new one: say the Telugu for to dance, then the Telugu for to celebrate.
\end{quote}

\subsection*{You'll want to know: \te{పలకరించు}}
\textbf{\te{పలకరించు}} — \emph{palakariñcu} — \textquotedblleft{}to greet\textquotedblright{}.

\begin{grammarlens}[title={What you've built}]
One more word for this chapter.
\end{grammarlens}

\subsection*{Your turn}
\begin{itemize}
\raggedright
  \item \emph{Say it:} \emph{palakariñcu}
  \item \emph{Say it:} \emph{palakariñcu}, once more
  \item \emph{Say it:} say \emph{palakariñcu}
  \item \emph{Recall:} say the Telugu for to dance, then the Telugu for to celebrate, then say \emph{palakariñcu} again
\end{itemize}

\begin{tcolorbox}[breakable,colback=teal!4,colframe=teal!35!black,title={Before you move on}]
What does \te{పలకరించు} mean? (\textquotedblleft{}To greet\textquotedblright{}.) Say it once more.
\end{tcolorbox}
\section[\texorpdfstring{\te{సిగ్గుపడు} (siggupaḍu)}{సిగ్గుపడు (siggupaḍu)}]{\te{సిగ్గుపడు} (siggupaḍu) — to feel shy, to be bashful}

\label{lesson:TE-C161-siggupadu}



\begin{quote}
Before the new one: say the Telugu for to celebrate, then the Telugu for to greet.
\end{quote}

\subsection*{You'll want to know: \te{సిగ్గుపడు}}
\textbf{\te{సిగ్గుపడు}} — \emph{siggupaḍu} — \textquotedblleft{}to feel shy, to be bashful\textquotedblright{}.

\begin{grammarlens}[title={What you've built}]
One more word for this chapter.
\end{grammarlens}

\subsection*{Your turn}
\begin{itemize}
\raggedright
  \item \emph{Say it:} \emph{siggupaḍu}
  \item \emph{Say it:} \emph{siggupaḍu}, once more
  \item \emph{Say it:} say \emph{siggupaḍu}
  \item \emph{Recall:} say the Telugu for to celebrate, then the Telugu for to greet, then say \emph{siggupaḍu} again
\end{itemize}

\begin{tcolorbox}[breakable,colback=teal!4,colframe=teal!35!black,title={Before you move on}]
What does \te{సిగ్గుపడు} mean? (\textquotedblleft{}To feel shy, to be bashful\textquotedblright{}.) Say it once more.
\end{tcolorbox}
\section[\texorpdfstring{\te{ఆశ్చర్యపోవు} (āścaryapōvu)}{ఆశ్చర్యపోవు (āścaryapōvu)}]{\te{ఆశ్చర్యపోవు} (āścaryapōvu) — to be surprised, to be amazed}

\label{lesson:TE-C161-ascaryapovu}



\begin{quote}
Before the new one: say the Telugu for to greet, then the Telugu for to feel shy.
\end{quote}

\subsection*{You'll want to know: \te{ఆశ్చర్యపోవు}}
\textbf{\te{ఆశ్చర్యపోవు}} — \emph{āścaryapōvu} — \textquotedblleft{}to be surprised, to be amazed\textquotedblright{}.

\begin{grammarlens}[title={What you've built}]
One more word for this chapter.
\end{grammarlens}

\subsection*{Your turn}
\begin{itemize}
\raggedright
  \item \emph{Say it:} \emph{āścaryapōvu}
  \item \emph{Say it:} \emph{āścaryapōvu}, once more
  \item \emph{Say it:} say \emph{āścaryapōvu}
  \item \emph{Recall:} say the Telugu for to greet, then the Telugu for to feel shy, then say \emph{āścaryapōvu} again
\end{itemize}

\begin{tcolorbox}[breakable,colback=teal!4,colframe=teal!35!black,title={Before you move on}]
What does \te{ఆశ్చర్యపోవు} mean? (\textquotedblleft{}To be surprised, to be amazed\textquotedblright{}.) Say it once more.
\end{tcolorbox}
